% 原图为 part1/visuals/nlp_pipeline.pdf 第 1 页中的 CoreNLP 系统示意图。
% 各分析器与标注对象双向连接；执行流单独向下，不新增分析器间的边。
\resizebox{.95\linewidth}{!}{%
\begin{tikzpicture}[
    font=\small,
    analyser/.style={draw,rounded corners=2pt,fill=blue!5,minimum width=4.4cm,minimum height=.54cm,align=center},
    object/.style={draw,rounded corners=5pt,fill=green!6,minimum width=2.2cm,minimum height=7.5cm,align=center},
    doc/.style={draw,fill=yellow!8,minimum width=1.55cm,minimum height=1cm,align=center},
    flow/.style={-{Stealth[length=2mm]},line width=.7pt}
]
\node[object] (obj) at (6.1,0) {标注对象};
\foreach \name/\yy/\zh/\code in {
    tok/3.35/词元化/tokenize,
    sent/2.4/分句/ssplit,
    pos/1.45/词性标注/pos,
    lemma/.5/形态分析/lemma,
    ner/-.45/命名实体识别/ner,
    parse/-1.4/句法分析/parse,
    coref/-2.35/共指消解/dcoref,
    other/-3.3/其他标注器/{gender, sentiment}}
{
    \node[analyser] (\name) at (1.8,\yy) {\zh};
    \node[font=\ttfamily\scriptsize,anchor=north,inner sep=1pt] at (1.8,\yy-.28) {(\code)};
    \draw[<->,>=Stealth,line width=.6pt] (\name.east) -- (\name.east -| obj.west);
}
\node[doc] (raw) at (9.1,1) {原始文本};
\node[doc] (annotated) at (9.1,-1.5) {已标注文本};
\draw[flow] (raw.west) -- (raw.west -| obj.east);
\draw[flow] (annotated.west -| obj.east) -- (annotated.west);
\draw[flow] (-.9,3.65) -- node[rotate=90,above,inner sep=4pt] {执行流} (-.9,-3.8);
\end{tikzpicture}%
}
